\documentclass[11pt]{article}
\usepackage[a4paper,margin=1in]{geometry}
\usepackage{amsmath,amssymb,amsthm,bm,mathtools}
\usepackage{physics}
\usepackage{booktabs}
\usepackage{hyperref}

\title{CV-DV LCHS Formulation Used in \texttt{formed\_qutip.py} and \texttt{formed\_bosonic.py}}
\author{Internal Derivation Note}
\date{\today}

\begin{document}
\maketitle

\section{Problem and Operator Model}
We solve linear ODE systems of the form
\begin{equation}
\frac{d}{dt}\bm{u}(t) = -A\,\bm{u}(t),
\qquad
\bm{u}(T)=e^{-AT}\bm{u}(0).
\end{equation}
In code, $A$ is provided by Pauli decomposition, split into two blocks:
\begin{equation}
L = \sum_j \ell_j P_j,
\qquad
H = \sum_k h_k Q_k,
\qquad
A=L+iH.
\end{equation}
(For the 1D heat-equation benchmark, $H=0$ and $L$ is the Laplacian block.)

\paragraph{Heat example used in the files.}
For $2$ DV qubits ($4\times 4$ system),
\begin{equation}
A = \frac{\alpha}{h_{\mathrm{grid}}^2}
\left(2\,II-IX-\frac12XX-\frac12YY\right).
\end{equation}

\section{CV-DV LCHS Representation}
The implemented hybrid map is
\begin{equation}
K(T)=\bra{\phi_r}\,e^{-iT\,\hat{x}\otimes A}\,\ket{\psi_{r,r',\beta}},
\label{eq:lchs-map}
\end{equation}
with
\begin{equation}
\ket{\phi_r}=S(r)\ket{0},
\qquad
\ket{\psi_{r,r',\beta}}=S(r')\sum_{n=0}^{N_c-1}C_n\ket{n}.
\end{equation}
The reference compares $K(T)$ to $e^{-AT}$ up to a global complex scale,
\begin{equation}
\eta^*=\arg\min_\eta\|K-\eta e^{-AT}\|_F,
\qquad
\epsilon_{\mathrm{map}}=\frac{\|K-\eta^*e^{-AT}\|_F}{\|\eta^*e^{-AT}\|_F}.
\end{equation}
For a given initial state $\bm{u}(0)$, the DV output vector
$\bm{u}_{\mathrm{cv}}=K(T)\bm{u}(0)$ is compared to the exact solution
$\bm{u}_{\mathrm{ref}}=e^{-AT}\bm{u}(0)$ via the \emph{state fidelity}
(shape-only, insensitive to the global scale~$\eta$):
\begin{equation}
F=\left|\frac{\langle\bm{u}_{\mathrm{cv}},\bm{u}_{\mathrm{ref}}\rangle}
{\|\bm{u}_{\mathrm{cv}}\|\;\|\bm{u}_{\mathrm{ref}}\|}\right|^2,
\end{equation}
and the \emph{post-selection probability}
\begin{equation}
p_{\mathrm{post}}=\|\bm{u}_{\mathrm{DV,post}}\|_2^2,
\end{equation}
where $\bm{u}_{\mathrm{DV,post}}$ is the DV register state after projecting
the CV mode onto the measurement bra $\bra{\phi_r}$.

\section{Coefficient Formula Used in Code}
\subsection{Kernel branch}
Following the near-optimal branch of An, Childs, and Lin~[1],
\begin{equation}
f_\beta(k)=\frac{e^{2^\beta}}{2\pi\,e^{(1+ik)^\beta}},
\qquad
g_\beta(k)=\frac{f_\beta(k)}{1-ik},
\qquad
0<\beta<1.
\end{equation}

\subsection{Corrected $\hbar=1$ expression}
Our QuTiP/bosonic implementation uses
\begin{equation}
\hat{x}=\frac{a+a^\dagger}{\sqrt{2}}
\quad (\hbar=1\text{ convention}).
\end{equation}
The corrected coefficient equation is
\begin{equation}
C_n \propto \int_{-\infty}^{\infty}
H_n\!\left(\frac{k}{\sigma'}\right)
\,g_\beta(k)
\,e^{-\gamma_{\hbar=1}k^2}
\,dk,
\qquad
\gamma_{\hbar=1}=\frac12\left(e^{-2r'}-e^{-2r}\right),
\label{eq:corrected-explicit}
\end{equation}
with $\sigma=e^r$, $\sigma'=e^{r'}$.

Equation \eqref{eq:corrected-explicit} is what \texttt{formed\_qutip.py} labels as the explicit-overlap backend.
The Gauss--Hermite compensated backend is the same formula under
\begin{equation}
k=\sqrt{2}\,\sigma'\xi,
\end{equation}
giving
\begin{equation}
C_n \propto \sum_i w_i
H_n(\sqrt{2}\xi_i)
\,g_\beta(\sqrt{2}\sigma'\xi_i)
\,\exp\!\left(\frac{\sigma'^2}{\sigma^2}\xi_i^2\right).
\label{eq:gh-comp}
\end{equation}

\section{Trotterized Gate Construction in \texttt{formed\_bosonic.py}}
For Pauli terms $A=\sum_m c_mP_m$, the first-order product formula~[6] is
\begin{equation}
e^{-iT\hat{x}\otimes A}
\approx
\left(\prod_m e^{-i\Delta t\,c_m\hat{x}\otimes P_m}\right)^{N_t},
\qquad
\Delta t=\frac{T}{N_t}.
\end{equation}
The code compiles each term by basis diagonalization of $P_m$ and parity mapping to one target qubit:
\begin{equation}
e^{-i\Delta t c_m\hat{x}\otimes P_m}
=
B_m^\dagger C_m^\dagger
\,e^{-i\Delta t c_m\hat{x}\otimes Z_t}
\,C_mB_m.
\end{equation}
Because
\begin{equation}
e^{-i\lambda\hat{x}}=D\!\left(-i\lambda/\sqrt{2}\right),
\end{equation}
the primitive gate is implemented with conditional displacement \texttt{cv\_c\_d}.

\section{CV State Preparation Modules and Theory}
\subsection{Injection (simulator baseline)}
\begin{equation}
\ket{\psi_{\mathrm{seed}}}=\sum_n C_n\ket{n}
\end{equation}
is directly loaded by \texttt{cv\_initialize}. This is not a physical pulse sequence, but gives a clean algorithmic baseline.

\subsection{Gate-based SNAP$+$D optimization}
We use the alternating ansatz
\begin{equation}
U_{\mathrm{SNAP+D}}=\prod_{\ell=1}^{L}D(\alpha_\ell)\,\mathrm{SNAP}(\bm{\theta}_\ell),
\end{equation}
optimized for
\begin{equation}
\max_{\{\alpha_\ell,\bm{\theta}_\ell\}}
\left|\bra{\psi_{\mathrm{seed}}}
U_{\mathrm{SNAP+D}}\ket{0}\right|^2.
\end{equation}
This follows the SNAP gate introduced by Heeres \emph{et al.}~[2] and the universal-oscillator-control framework (including SNAP-like number-selective phases plus displacements) of Krastanov \emph{et al.}~[12]. We view this gate set through the hybrid CV--DV instruction-set / abstract-machine perspective summarized by Liu \emph{et al.}~[3].

\subsection{Optimization-free Givens decomposition (Law--Eberly protocol)}
\label{sec:givens}
The Givens method prepares $\ket{\psi_{\mathrm{seed}}}=\sum_{n=0}^{N-1}C_n\ket{n}$
from the vacuum $\ket{0}$ using $N-1$ analytically determined rotation steps,
with \emph{no numerical optimization}.

\paragraph{Givens rotation.}
Define the two-level rotation acting on adjacent Fock levels $\ket{n}$, $\ket{n+1}$:
\begin{equation}
G(n,n{+}1;\theta,\varphi)
=\begin{pmatrix}\cos\theta & -e^{-i\varphi}\sin\theta\\
e^{i\varphi}\sin\theta & \cos\theta\end{pmatrix}
\oplus I_{\perp},
\label{eq:givens-rotation}
\end{equation}
where the $2\times2$ block acts on the $\{\ket{n},\ket{n+1}\}$ subspace
and $I_\perp$ is the identity on all other Fock levels.
Any $N$-component target state can be factored as
\begin{equation}
\ket{\psi_{\mathrm{seed}}}
=G(N{-}2,N{-}1;\theta_{N-1},\varphi_{N-1})\cdots G(1,2;\theta_2,\varphi_2)
\,G(0,1;\theta_1,\varphi_1)\,\ket{0},
\label{eq:givens-decomposition}
\end{equation}
where the angles $\{(\theta_k,\varphi_k)\}_{k=1}^{N-1}$ are computed
recursively from the coefficients $C_n$ by a top-down elimination
(see \texttt{formed\_stateprep.py}, function \texttt{givens\_decomposition}).
This decomposition is \emph{exact}: the prepared state equals the target
to machine precision, with fidelity $|\braket{\psi_{\mathrm{prep}}}{\psi_{\mathrm{seed}}}|^2=1$.

\paragraph{Law--Eberly physical realization.}
On a coupled qubit--oscillator platform, a Law--Eberly synthesis alternates
(i) single-qubit rotations that set relative amplitudes/phases and
(ii) a calibrated qubit--oscillator \emph{exchange} interaction that transfers
population between adjacent Fock levels.
The exchange dynamics are generated by the Jaynes--Cummings Hamiltonian~[4]
\begin{equation}
H_{\mathrm{JC}} = \hbar\Omega(t)\bigl(\sigma_- a^\dagger + \sigma_+ a\bigr),
\label{eq:jc-hamiltonian}
\end{equation}
which couples $\ket{e,n}\leftrightarrow\ket{g,n{+}1}$ at the $\sqrt{n+1}$-dependent rate
$\Omega_{n}=\Omega_0\sqrt{n+1}$.
Importantly, a \emph{single} resonant JC ``swap'' pulse is not photon-number selective by itself; the
effective selectivity required by the Givens factorization is achieved in practice by
hardware calibration and, when available, photon-number-resolved qubit control in the dispersive regime
(Section~\ref{sec:jc-selectivity}).
A representative experimental demonstration of deterministic resonator-state synthesis using such calibrated
sequences is Hofheinz~\emph{et al.}~[11].

\paragraph{Theorem (Law and Eberly~[5], Eqs.~5--7).}
\emph{Any Fock superposition $\sum_{n=0}^{N-1}C_n\ket{n}$ can be prepared
from $\ket{0,g}$ in exactly $N-1$ qubit-rotation + selective-JC steps.}
The total gate cost is:
\begin{equation}
\text{Resources:}\quad
(N_{\mathrm{active}}-1)\text{ JC pulses}
+\,(N_{\mathrm{active}}-1)\text{ qubit rotations},
\label{eq:givens-resources}
\end{equation}
where $N_{\mathrm{active}}=|\{n:C_n\ne 0\}|$ is the number of occupied Fock levels.

\paragraph{Simulator realization.}
The bosonic-qiskit library provides a \emph{global} JC gate
$U_{\mathrm{JC}}(\theta,\phi)=\exp\bigl(-i\theta(\sigma_- a^\dagger e^{i\phi}+\mathrm{h.c.})\bigr)$,
which couples \emph{all} Fock manifolds simultaneously with
$\sqrt{n+1}$-dependent Rabi frequencies.
A single \texttt{cv\_jc} call therefore cannot implement a
number-selective Givens rotation without crosstalk to other manifolds.
In the current implementation,
the analytically computed state $\ket{\psi_{\mathrm{seed}}}$ is loaded
via \texttt{cv\_initialize} (simulator injection), and the Givens gate
count~\eqref{eq:givens-resources} is reported for physical resource estimation.
On hardware with frequency-selective JC (e.g., superconducting transmon--cavity),
the same angles would translate directly into pulse parameters.

\section{Documented Failure Cases and Lessons}
\subsection{Failed CV prep by inconsistency aggregation}
A previous internal strategy aggregated inconsistencies over map error, coefficient-backend mismatch, and truncation tail mass. A representative objective is
\begin{equation}
\mathcal{J}(\beta)=
\lambda_1\,\epsilon_{\mathrm{map}}(\beta)
+\lambda_2\,\epsilon_{\mathrm{backend}}(\beta)
+\lambda_3\,\tau_{\mathrm{tail}}(\beta),
\end{equation}
where
\begin{equation}
\epsilon_{\mathrm{backend}}=
\frac{\|\bm{C}^{(\mathrm{exp})}-\eta\bm{C}^{(\mathrm{gh})}\|_2}{\|\eta\bm{C}^{(\mathrm{gh})}\|_2},
\qquad
\tau_{\mathrm{tail}}=\sum_{n=N_c-4}^{N_c-1}|C_n|^2.
\end{equation}
The failure mode was not one single bug, but amplified disagreement when both backend mismatch and truncation pressure were high.

Empirically we often observed better aggregate consistency at lower $\beta$ (within admissible range), because the phase-growth term in
\begin{equation}
(1+ik)^\beta=|1+ik|^\beta e^{i\beta\arg(1+ik)}
\end{equation}
becomes less oscillatory as $\beta$ decreases, which can reduce destructive quadrature cancellation in finite precision.
This is a numerical-stability trend, not a universal theorem; final choice still depends on $r,r',N_c,N_{\mathrm{quad}}$.

\subsection{Current numerical objective in \texttt{formed\_sweep.py}}
The current sweep/optimization workflow treats final-state fidelity and
post-selection probability as \emph{joint} objectives.
For each valid run, we compute
\begin{equation}
F=\text{state fidelity},\qquad
p_{\mathrm{post}}=\|\bm{u}_{\mathrm{DV,post}}\|_2^2.
\end{equation}
Ranking is multi-objective:
\begin{enumerate}
\item Pareto-front rank on $(F,p_{\mathrm{post}})$ (maximize both).
\item Deterministic tie-break by
\begin{equation}
s_{\mathrm{mo}}=\sqrt{F\,p_{\mathrm{post}}}.
\end{equation}
\end{enumerate}
Local continuous refinement minimizes
\begin{equation}
\mathcal{L}_{\mathrm{local}}=1-s_{\mathrm{mo}}.
\end{equation}
In the current default sweep box, parameters are explored over
$r\in[0.5,5]$, $r'\in[0.01,4]$, and $\beta\in[0.3,0.99]$, with the
physical constraint $r'<r$.

\subsection{Number selectivity, pulse-level control, and why we use injection}
\label{sec:jc-selectivity}
The Givens decomposition (Section~\ref{sec:givens}) assumes access to
\emph{selective} rotations that act only on a chosen adjacent Fock subspace
$\{\ket{n},\ket{n{+}1}\}$.
On real hybrid qubit--oscillator hardware, this selectivity is obtained through
a combination of (i) calibrated exchange dynamics (JC swaps) and
(ii) photon-number-resolved qubit control in the dispersive regime.
Gate-level CV simulators generally expose the ideal JC unitary, but they do not
model the pulse-level selectivity mechanisms as primitive gates; this motivates
the simulation strategy used in \texttt{formed\_bosonic.py}.

\paragraph{Dispersive Hamiltonian and photon-number splitting.}
In circuit QED, a qubit coupled to a cavity is in the \emph{dispersive regime} when
the detuning $\Delta=\omega_q-\omega_c$ is large compared to the coupling $g$,
i.e., $|\Delta|\gg g$~[13,14].
To leading order, the effective Hamiltonian (up to constant shifts) can be written as~[14]
\begin{equation}
H_{\mathrm{disp}}
=\hbar\omega_c\,a^\dagger a
+\frac{\hbar\omega_q}{2}\sigma_z
-\hbar\chi\,a^\dagger a\;\sigma_z,
\qquad
\chi \approx \frac{g^2}{\Delta},
\label{eq:dispersive-H}
\end{equation}
so the qubit transition frequency becomes photon-number dependent,
\begin{equation}
\omega_q^{(n)}=\omega_q - 2\chi\,n.
\label{eq:number-dependent-freq}
\end{equation}
This photon-number splitting was directly observed in superconducting circuits by
Schuster~\emph{et al.}~[15].

\paragraph{What is actually ``number selective'' in practice.}
Photon-number dependence enables \emph{number-selective qubit rotations}:
a narrow-band qubit drive centered at $\omega_q^{(n)}$ will rotate the qubit
\emph{only} when the cavity contains $n$ photons, provided the drive bandwidth is
much smaller than the splitting, i.e., $\sigma_\omega\ll 2|\chi|$~[14,15].
These number-selective qubit controls are the physical basis for SNAP-type operations~[2]
and are also used as building blocks in deterministic resonator-state synthesis
protocols such as the Law--Eberly construction and its experimental realization by
Hofheinz~\emph{et al.}~[11].

In contrast, resonant JC exchange (Section~\ref{sec:givens}) is typically implemented by tuning the qubit close
to resonance with the cavity for a calibrated time.  Such an exchange pulse is \emph{global} in photon number,
so selectivity is achieved by the overall calibrated sequence (and/or by interleaving dispersive, number-selective
qubit operations), not by an ideal manifold-isolated JC gate.

\paragraph{Why JC-only gate-level simulation is insufficient for dispersive selectivity.}
Gate-level simulators (bosonic-qiskit~[16], hybridlane~[17]) expose useful idealized selective primitives
(e.g., SNAP-type operations), but their JC interaction is still an ideal global unitary
\begin{equation}
U_{\mathrm{JC}}(\theta,\phi)
=\exp\!\bigl(-i\theta(\sigma_-\,a^\dagger e^{i\phi}+\sigma_+\,a\,e^{-i\phi})\bigr),
\label{eq:global-jc}
\end{equation}
which couples \emph{all} Fock manifolds simultaneously at $\sqrt{n+1}$-dependent rates.
A pulse calibrated for one manifold therefore produces coherent ``crosstalk'' on all others, so a direct
implementation of Eq.~\eqref{eq:givens-decomposition} using repeated global \texttt{cv\_jc} calls does not realize the intended
selective Givens rotations.

To model the \emph{physical} origin of dispersive selectivity from first principles, one would need a pulse-level model with explicit
dispersive parameters (e.g., time-dependent Hamiltonian integration).  Using number-selective qubit rotations / SNAP as gate-level primitives
is still valid as an \emph{effective abstraction}, but then selectivity is assumed rather than emergent from the JC pulse dynamics. The latter viewpoint is emphasized in
the hybrid CV--DV ISA / abstract-machine treatment of Liu~\emph{et al.}~[3].

\paragraph{Justification for injection in simulation.}
Given this gap, \texttt{formed\_bosonic.py} uses \texttt{cv\_initialize} (direct Fock-amplitude injection) for the
Givens/Law--Eberly branch and separately reports the analytic pulse count from the Givens decomposition.
This follows a standard separation-of-concerns strategy:
\begin{enumerate}
\item \textbf{Correctness.} Injection loads the exact target state
  $\ket{\psi_{\mathrm{seed}}}=\sum_n C_n\ket{n}$, isolating LCHS approximation errors
  (kernel truncation, Trotter error, and post-selection statistics) from device-dependent control errors.
\item \textbf{Physical feasibility.} Deterministic resonator-state synthesis with calibrated qubit--resonator control
  has been demonstrated experimentally~[11], and SNAP-based control of oscillator phases is experimentally established~[2].
\item \textbf{Resource accounting.} The Givens factorization yields an exact \emph{count} of required elementary steps
  (Eq.~\eqref{eq:givens-resources}), while SNAP$+$D provides a gate-level (optimization-based) alternative grounded in
  established universality results~[12].
\item \textbf{Precedent.} Simulations of bosonic-encoded protocols commonly assume ideal preparation of nontrivial oscillator states
  (e.g., GKP~[19], cat codes~[20]) while reporting gate/pulse counts separately.
\end{enumerate}
\section{Practical Recommendations Used by the Two Formed Files}
\begin{itemize}
\item Keep one convention end-to-end ($\hat{x}=(a+a^\dagger)/\sqrt{2}$ in both QuTiP and bosonic mapping).
\item Check explicit-overlap vs GH-comp coefficients before running expensive gate-level circuits.
\item \textbf{CV state preparation strategy} (see Table~\ref{tab:stateprep}):
  \begin{itemize}
  \item \emph{Injection} (\texttt{cv\_initialize}): exact, simulator-only.
    Use as the theoretical baseline against which gate-based methods are compared.
  \item \emph{SNAP$+$D}: gate-compiled, optimization-based.
    Best simulator fidelity ($>99.99\%$ at depth~4).
    Suitable for circuit-level benchmarking.
  \item \emph{Givens}: optimization-free, analytically exact angles.
    Reports the physical JC gate count; uses injection in the simulator
    due to the global-JC selectivity gap (Section~\ref{sec:jc-selectivity}).
    Suitable for hardware resource estimation.
  \end{itemize}
\item Track both map-level error ($\epsilon_{\mathrm{map}}$) and vector-level fidelity after post-selection.
\item In sweep optimization, treat fidelity and post-selection probability as
  joint objectives (Pareto-first ranking with scalar tie-break
  $s_{\mathrm{mo}}=\sqrt{F\,p_{\mathrm{post}}}$).
\item Increase $N_c$ and Fock cutoff together when tail mass is not small.
\end{itemize}

\begin{table}[h]
\centering
\caption{CV state preparation methods implemented in \texttt{formed\_bosonic.py}.}
\label{tab:stateprep}
\begin{tabular}{lcccl}
\toprule
Method & Optimization & Gate-based & Fidelity & Reference \\
\midrule
Injection   & --- & No  & $1.0$ (exact) & --- \\
SNAP$+$D    & L-BFGS-B & Yes & $>0.9999$ (depth~4) & Heeres~[2], Krastanov~[12] \\
Givens      & None     & Yes$^*$ & $1.0$ (analytic) & Law--Eberly~[5], Hofheinz~[11] \\
\bottomrule
\multicolumn{5}{l}{\footnotesize $^*$Gate-compiled on hardware with number-selective JC; simulator uses injection.}
\end{tabular}
\end{table}

\section{Library Convention Notes (with Links)}
\paragraph{QuTiP.}
QuTiP directly documents the single-mode position operator as
\[
x=\frac{1}{\sqrt{2}}(a+a^\dagger).
\]
See:
\begin{itemize}
\item \url{https://qutip.readthedocs.io/en/latest/apidoc/quantumobject.html#qutip.core.operators.position}
\end{itemize}
For CV helper functions, QuTiP also documents the scaling relation
\[
\hbar = \frac{2}{g^2},
\]
with default $g=\sqrt{2}$, giving default $\hbar=1$. See:
\begin{itemize}
\item \url{https://qutip.readthedocs.io/en/stable/_modules/qutip/continuous_variables.html}
\end{itemize}

\paragraph{bosonic-qiskit.}
The bosonic-qiskit gate API does not expose a user-level \texttt{hbar} argument in
\texttt{cv\_d} (or other CV primitives), so convention control is handled by analytic mapping
outside the API. See:
\begin{itemize}
\item \url{https://c2qa.github.io/bosonic-qiskit/c2qa.circuit.CVCircuit.html#c2qa.circuit.CVCircuit.cv_d}
\end{itemize}
Its displacement operator implementation is
\[
D(\alpha)=\exp(\alpha a^\dagger-\alpha^* a),
\]
from source:
\begin{itemize}
\item \url{https://raw.githubusercontent.com/C2QA/bosonic-qiskit/cee0d932e5dd7855dd1eeefc33a488c4c35e9ef0/src/c2qa/operators.py}
\end{itemize}
Therefore, to match a target convention (e.g., $\hat{x}=(a+a^\dagger)/\sqrt{2}$), we must
apply the correct parameter conversion in circuit construction.

\section{Photon Loss Noise Simulation}
\label{sec:noise}

\subsection{Photon loss channel}
The dominant decoherence mechanism for a bosonic mode coupled to a zero-temperature
bath is single-photon loss at rate~$\kappa$ (units of~s$^{-1}$).
Over a time interval~$\Delta t$, the channel acts on the CV mode density matrix as~[7,21]
\begin{equation}
\mathcal{E}_\kappa(\rho)=\sum_{n=0}^{\infty}K_n\,\rho\,K_n^\dagger,
\label{eq:photon-loss-channel}
\end{equation}
with Kraus operators
\begin{equation}
K_n=\sqrt{\frac{(1-e^{-\kappa\Delta t})^n}{n!}}\;
e^{-\kappa\Delta t\,\hat{N}/2}\;a^n,
\qquad
\hat{N}=a^\dagger a.
\label{eq:kraus-photon-loss}
\end{equation}
For $\kappa=0$ the only surviving Kraus operator is $K_0=I$, recovering the noiseless channel.

In the bosonic-qiskit implementation~[16], each CV gate carries a default
duration of $\Delta t_{\mathrm{gate}}=100\;\text{ns}$, and the photon loss channel
is applied after every gate via the \texttt{PhotonLossNoisePass} transpiler pass.
For a circuit with $N_{\mathrm{gate}}$ bosonic operations, the total CV-mode exposure
time is
\begin{equation}
T_{\mathrm{exp}}=N_{\mathrm{gate}}\,\Delta t_{\mathrm{gate}}.
\label{eq:exposure-time}
\end{equation}
In the current 2-qubit heat-equation circuit ($N_t=10$ Trotter steps,
4 Pauli terms with $\sim\!5$ bosonic ops each, plus state-preparation and
final squeeze), $N_{\mathrm{gate}}\approx 52$, giving
$T_{\mathrm{exp}}\approx 5.2\;\mu\text{s}$.

\subsection{Density-matrix simulation path}
\label{sec:dm-path}
Because the Kraus channel~\eqref{eq:photon-loss-channel} produces a
mixed state, the simulation must evolve a density matrix rather than a
statevector.
In \texttt{formed\_bosonic.py} this is achieved by:
\begin{enumerate}
\item Appending a \texttt{save\_density\_matrix()} instruction to the
  circuit (in place of \texttt{save\_statevector}).
\item Calling \texttt{cv\_util.simulate()} with
  \texttt{add\_save\_statevector=False} and passing the noise pass
  via the \texttt{noise\_passes} argument.
  The simulator automatically selects the Aer \texttt{density\_matrix}
  method when Kraus channels are present.
\end{enumerate}
The output is the full joint density matrix
$\rho\in\mathbb{C}^{D\times D}$ with $D=N_{\mathrm{fock}}\cdot 2^{n_q}$.

\subsection{Post-selected DV density matrix}
After the circuit applies the final squeeze $S^\dagger(r)$ (already part
of \texttt{build\_bosonic\_circuit}), post-selection on the CV vacuum
$\ket{0}$ is implemented by extracting the sub-block of~$\rho$
corresponding to Fock index~0:
\begin{equation}
(\rho_{\mathrm{DV}})_{ij}
=\bra{0,i}\rho\ket{0,j},
\qquad
i,j\in\{0,\dots,2^{n_q}-1\}.
\label{eq:postselect-dm}
\end{equation}
The matrix $\rho_{\mathrm{DV}}$ is Hermitian positive-semidefinite but generally
\emph{sub-normalized}: its trace equals the post-selection probability,
\begin{equation}
p_{\mathrm{post}}=\operatorname{Tr}(\rho_{\mathrm{DV}}).
\label{eq:post-prob-dm}
\end{equation}

\paragraph{Layout-aware indexing.}
The bosonic-qiskit simulator uses one of two memory layouts for the
joint Hilbert space $\mathcal{H}_{\mathrm{fock}}\otimes\mathcal{H}_{\mathrm{DV}}$:
\begin{itemize}
\item \textbf{fock\_major}:
  index $= f\cdot 2^{n_q} + d$, so the Fock-0 block occupies the
  top-left $2^{n_q}\times 2^{n_q}$ corner.
\item \textbf{qubit\_major}:
  index $= d\cdot N_{\mathrm{fock}} + f$, so the Fock-0 entries sit at
  stride-$N_{\mathrm{fock}}$ positions $\{0,N_{\mathrm{fock}},2N_{\mathrm{fock}},\dots\}$.
\end{itemize}
In both cases, a bit-reversal permutation is applied to the extracted
$2^{n_q}\times 2^{n_q}$ block to convert from qiskit's little-endian
qubit ordering to physics convention.
The layout is detected once per parameter set via a small probe
simulation (\texttt{detect\_statevector\_layout}).

\subsection{Mixed-state fidelity}
For a sub-normalized post-selected density matrix $\rho_{\mathrm{DV}}$ and
a pure reference vector $\bm{u}_{\mathrm{ref}}$, the fidelity is
\begin{equation}
F=\frac{\bm{u}_{\mathrm{ref}}^\dagger\,\rho_{\mathrm{DV}}\,\bm{u}_{\mathrm{ref}}}
{\operatorname{Tr}(\rho_{\mathrm{DV}})\;\|\bm{u}_{\mathrm{ref}}\|^2}.
\label{eq:fidelity-dm}
\end{equation}
This reduces to the standard pure-state fidelity
$|\!\braket{\bm{u}_{\mathrm{cv}}}{\bm{u}_{\mathrm{ref}}}\!|^2/
(\|\bm{u}_{\mathrm{cv}}\|^2\|\bm{u}_{\mathrm{ref}}\|^2)$
when $\rho_{\mathrm{DV}}$ is rank-1.
A zero-trace guard returns $F=0$ whenever
$\operatorname{Tr}(\rho_{\mathrm{DV}})<10^{-15}$.

\subsection{Purity}
The purity of the (renormalized) post-selected state quantifies
the degree of mixing introduced by the loss channel:
\begin{equation}
\gamma
=\frac{\operatorname{Tr}(\rho_{\mathrm{DV}}^2)}
      {\bigl[\operatorname{Tr}(\rho_{\mathrm{DV}})\bigr]^2}.
\label{eq:purity}
\end{equation}
For a pure state $\gamma=1$; photon loss generically decreases purity
but the relationship is not guaranteed monotonic in~$\kappa$ because
post-selection on Fock~$\ket{0}$ preferentially retains population that
has already decayed to vacuum.

\subsection{Noise sweep methodology}
The $\kappa$-sweep in \texttt{formed\_bosonic.py} (gated by
\texttt{RUN\_NOISE\_SWEEP}) evaluates the injection-based LCHS circuit
at 20 logarithmically spaced loss rates,
\begin{equation}
\kappa\in\bigl\{10^3,\;\dots,\;10^6\bigr\}\;\text{s}^{-1},
\label{eq:kappa-range}
\end{equation}
chosen so that the product $\kappa\,T_{\mathrm{exp}}$ spans the
range $\sim\!5\times 10^{-3}$ (negligible loss) to $\sim\!5$
(strong loss) for $T_{\mathrm{exp}}\approx 5.2\;\mu\text{s}$.

For each value of~$\kappa$, the sweep records:
\begin{itemize}
\item Fidelity vs.\ classical reference
  $F_{\mathrm{cl}}$ (Eq.~\ref{eq:fidelity-dm} with
  $\bm{u}_{\mathrm{ref}}=e^{-AT}\bm{u}(0)$).
\item Fidelity vs.\ QuTiP reference
  $F_{\mathrm{qt}}$ (Eq.~\ref{eq:fidelity-dm} with $\bm{u}_{\mathrm{ref}}$ from the
  noiseless QuTiP simulation).
\item Post-selection probability $p_{\mathrm{post}}$
  (Eq.~\ref{eq:post-prob-dm}).
\item Purity $\gamma$ (Eq.~\ref{eq:purity}).
\end{itemize}
Results are saved to \texttt{results\_formed/noise/noise\_sweep.csv}.

\section*{References}
\begin{enumerate}
\item D.~An, A.~M.~Childs, and L.~Lin, ``Quantum Algorithm for Linear Non-unitary Dynamics with Near-Optimal Dependence on All Parameters,''
\emph{Communications in Mathematical Physics} \textbf{407}(1) (2025). DOI: \url{https://doi.org/10.1007/s00220-025-05509-w}. Also arXiv:2312.03916.
\item R.~W.~Heeres \emph{et al.}, ``Cavity State Manipulation Using Photon-Number Selective Phase Gates,'' \emph{Phys.\ Rev.\ Lett.}\ \textbf{115}, 137002 (2015).
\item Y.~Liu \emph{et al.}, ``Hybrid Oscillator--Qubit Quantum Processors: Instruction Set Architectures, Abstract Machine Models, and Applications,'' arXiv:2407.10381.
\item E.~T.~Jaynes and F.~W.~Cummings, ``Comparison of quantum and semiclassical radiation theories with application to the beam maser,'' \emph{Proc.\ IEEE} \textbf{51}, 89 (1963). \url{https://doi.org/10.1109/PROC.1963.1664}
\item C.~K.~Law and J.~H.~Eberly, ``Arbitrary control of a quantum electromagnetic field,'' \emph{Phys.\ Rev.\ Lett.}\ \textbf{76}, 1055 (1996). \url{https://doi.org/10.1103/PhysRevLett.76.1055}
\item G.~Strang, ``On the construction and comparison of difference schemes,'' \emph{SIAM J.\ Numer.\ Anal.}\ \textbf{5}, 506 (1968). \url{https://doi.org/10.1137/0705041}
\item C.~Weedbrook \emph{et al.}, ``Gaussian quantum information,'' \emph{Rev.\ Mod.\ Phys.}\ \textbf{84}, 621 (2012). arXiv:1110.3234.
\item A.~Serafini, \emph{Quantum Continuous Variables: A Primer of Theoretical Methods} (CRC Press, 2017), Ch.~2.
\item S.~L.~Braunstein and P.~van~Loock, ``Quantum information with continuous variables,'' \emph{Rev.\ Mod.\ Phys.}\ \textbf{77}, 513 (2005). arXiv:quant-ph/0410100.
\item C.~C.~Gerry and P.~L.~Knight, \emph{Introductory Quantum Optics} (Cambridge University Press, 2004), Sec.~2.1.
\item M.~Hofheinz \emph{et al.}, ``Synthesizing arbitrary quantum states in a superconducting resonator,'' \emph{Nature} \textbf{459}, 546 (2009). \url{https://doi.org/10.1038/nature08005}
\item S.~Krastanov \emph{et al.}, ``Universal control of an oscillator with dispersive coupling to a qubit,'' \emph{Phys.\ Rev.\ A}\ \textbf{92}, 040303(R) (2015). \url{https://doi.org/10.1103/PhysRevA.92.040303}
\item A.~Blais \emph{et al.}, ``Cavity quantum electrodynamics for superconducting electrical circuits: An architecture for quantum computation,'' \emph{Phys.\ Rev.\ A}\ \textbf{69}, 062320 (2004). \url{https://doi.org/10.1103/PhysRevA.69.062320}
\item A.~Blais \emph{et al.}, ``Circuit quantum electrodynamics,'' \emph{Rev.\ Mod.\ Phys.}\ \textbf{93}, 025005 (2021). \url{https://doi.org/10.1103/RevModPhys.93.025005}
\item D.~I.~Schuster \emph{et al.}, ``Resolving photon number states in a superconducting circuit,'' \emph{Nature} \textbf{445}, 515 (2007). \url{https://doi.org/10.1038/nature05461}
\item T.~J.~Stavenger \emph{et al.}, ``Bosonic Qiskit,'' arXiv:2209.11153.
\item \emph{hybridlane}: a PennyLane plugin for hybrid CV--DV simulation (Pacific Northwest National Laboratory), \url{https://github.com/pnnl/hybridlane}.
\item P.~Reinhold \emph{et al.}, ``Error-corrected gates on an encoded qubit,'' \emph{Nat.\ Phys.}\ \textbf{16}, 822 (2020); see also P.~Eickbusch \emph{et al.}, ``Fast universal control of an oscillator with weak dispersive coupling to a qubit,'' \emph{Nat.\ Phys.}\ \textbf{18}, 1464 (2022).
\item D.~Gottesman, A.~Kitaev, and J.~Preskill, ``Encoding a qubit in an oscillator,'' \emph{Phys.\ Rev.\ A}\ \textbf{64}, 012310 (2001). \url{https://doi.org/10.1103/PhysRevA.64.012310}
\item M.~Mirrahimi \emph{et al.}, ``Dynamically protected cat-qubits: a new paradigm for universal quantum computation,'' \emph{New J.\ Phys.}\ \textbf{16}, 045014 (2014). \url{https://doi.org/10.1088/1367-2630/16/4/045014}
\item H.~J.~Carmichael, \emph{Statistical Methods in Quantum Optics 1: Master Equations and Fokker--Planck Equations} (Springer, 1999), Ch.~1. \url{https://doi.org/10.1007/978-3-662-03875-8}
\end{enumerate}

\end{document}
